\documentclass[../../../include/open-logic-section]{subfiles}

\begin{document}
\olfileid{sth}{card-arithmetic}{simp}

\olsection{简化加法与乘法}

事实证明，超限基数的加法和乘法是\emph{极其}简单的。这是因为基数是（某种）序数，因而是良序的，所以能够以某种特定的方式来操作。然而，证明这一点可\emph{并不}容易。首先，我们需要一个巧妙的定义：

\begin{defn}
我们在序数对上定义\emph{典范排序} $\canonord$，规定 $\tuple{\alpha_1, \alpha_2} \canonord
\tuple{\beta_1, \beta_2}$ 当且仅当满足下列条件之一：
\begin{enumerate}
	\item $\max(\alpha_1, \alpha_2) < \max(\beta_1, \beta_2)$；或
	\item $\max(\alpha_1, \alpha_2) = \max(\beta_1, \beta_2)$ 且 $\alpha_1 < \beta_1$；或
	\item $\max(\alpha_1, \alpha_2) = \max(\beta_1, \beta_2)$ 且 $\alpha_1 = \beta_1$ 且 $\alpha_2 < \beta_2$。
\end{enumerate}
\end{defn}

\begin{lem}
对于任意序数 $\alpha$，$\tuple{\alpha \times \alpha, \canonord}$ 是一个良序。
\end{lem}

\begin{proof}
显然 $\canonord$ 在 $\alpha \times \alpha$ 上是连通的。假设 $\tuple{\alpha_1, \alpha_2}$ 和 $\tuple{\beta_1, \beta_2}$ 互不 $\canonord$-小于对方。那么 $\max(\alpha_1,
\alpha_2) = \max(\beta_1, \beta_2)$，$\alpha_1 = \beta_1$ 且 $\alpha_2 = \beta_2$，因此 $\tuple{\alpha_1, \alpha_2} =
\tuple{\beta_1,  \beta_2}$。

为证明它是良序，设 $X \subseteq \alpha\times\alpha$ 非空。由于 $\alpha$ 是序数，某个 $\delta$ 是 $\Setabs{\max(\gamma_1, \gamma_2)}{\tuple{\gamma_1,
\gamma_2} \in X}$ 的最小元。现在，从 $\Setabs{\tuple{\gamma_1,\gamma_2} \in X}{\max(\gamma_1, \gamma_2) =
\delta}$ 中剔除第一坐标不是最小的所有序对；在剩下的序对中，第二坐标最小的那个就是 $X$ 的 $\canonord$-最小元。
\end{proof}
\noindent
现在，我们做一个微小且简单的观察：

\begin{prop}\ollabel{simplecardproduct}
如果 $\cardeq{\alpha}{\beta}$，那么 $\cardeq{\alpha \times \alpha}{\beta
\times \beta}$。
\end{prop}

\begin{proof}
只需令 $f \colon \alpha \to \beta$ 诱导出 $\tuple{\gamma_1,
\gamma_2} \mapsto \tuple{f(\gamma_1), f(\gamma_2)}$。
\end{proof}
\noindent
现在我们就要把这些都用上，来证明一个关键的引理：
\begin{lem}\ollabel{alphatimesalpha}
对任意无限序数 $\alpha$，有 $\cardeq{\alpha}{\alpha \times \alpha}$。
\end{lem}

\begin{proof}
为归谬，令 $\alpha$ 是使该结论不成立的最小无限序数。\olref[sfr][siz][zigzag]{natsquaredenumerable} 表明 $\cardeq{\omega}{\omega\times\omega}$，所以 $\omega \in \alpha$。此外，$\alpha$ 是一个基数：为归谬，假设并非如此；那么 $\card{\alpha} \in \alpha$，于是由假设，有 $\cardeq{\card{\alpha}}{\card{\alpha} \times \card{\alpha}}$；而根据定义有 $\cardeq{\card{\alpha}}{\alpha}$；因此由 \olref{simplecardproduct} 可得 $\cardeq{\alpha}{\alpha\times\alpha}$。

现在，对每个 $\tuple{\gamma_1, \gamma_2} \in \alpha \times \alpha$，考虑其前段：
\begin{align*}
	\text{Seg}(\gamma_1, \gamma_2) &= \Setabs{\tuple{\delta_1, \delta_2} \in \alpha \times \alpha}{\tuple{\delta_1, \delta_2} \canonord \tuple{\gamma_1, \gamma_2}}
\end{align*}
令 $\gamma = \max(\gamma_1, \gamma_2)$，注意 $\tuple{\gamma_1, \gamma_2} \canonord \tuple{\gamma+1, \gamma + 1}$。那么，当 $\gamma$ 无限时，观察：
\begin{align*}
	\text{Seg}(\gamma_1, \gamma_2) &
	\precsim ((\gamma \ordplus 1)\ordtimes (\gamma \ordplus 1))\\
	&\approx (\gamma \ordtimes \gamma)
	\text{，根据 \olref[ord-arithmetic][using-addition]{ordinfinitycharacter} 和
	\olref{simplecardproduct}}\\
	&\approx \gamma \text{，由归纳假设}\\
	& \prec \alpha\text{，因为 $\alpha$ 是一个基数}
\end{align*}
所以 $\ordtype{\alpha\times \alpha, \canonord} \leq \alpha$，从而 $\cardle{\alpha \times \alpha}{\alpha}$。既然显然有 $\cardle{\alpha}{\alpha \times \alpha}$，那么根据 Schr\"oder-Bernstein 定理即得结论。
\end{proof}

最后，我们得出简化结果：

\begin{thm}\ollabel{cardplustimesmax}
如果 $\cardfont{a}, \cardfont{b}$ 是无限基数，那么：
\[
	\cardfont{a}
\cardtimes \cardfont{b} = \cardfont{a} \cardplus \cardfont{b} =
\text{最大值}(\cardfont{a}, \cardfont{b}).
\]
\end{thm}

\begin{proof}
不失一般性，假设 $\cardfont{a} = \max(\cardfont{a},
\cardfont{b})$。那么由 \olref{alphatimesalpha}，有 $\cardfont{a}\cardtimes\cardfont{a} = \cardfont{a} \leq \cardfont{a}
\cardplus \cardfont{b} \leq \cardfont{a} \cardplus \cardfont{a} \leq
\cardfont{a} \cardtimes \cardfont{a}$。 \end{proof}\noindent 类似地，若 $\cardfont{a}$ 是无限基数，则由大小不超过 $\cardfont{a}$ 的集合构成的大小为 $\cardfont{a}$ 的并集，其大小也不超过 $\cardfont{a}$：

\begin{prop}\ollabel{kappaunionkappasize}
设 $\cardfont{a}$ 是一个无限基数。对每个序数 $\beta
\in \cardfont{a}$，令 $X_\beta$ 是一个满足 $\card{X_\beta} \leq
\cardfont{a}$ 的集合。那么 $\card{\bigcup_{\beta \in \cardfont{a}} X_\beta}
\leq \cardfont{a}$。
\end{prop}

\begin{proof}
对每个 $\beta \in \cardfont{a}$，取定一个单射 $f_\beta \colon
X_\beta \to \cardfont{a}$。\footnote{这些单射是如何“取定”的？参见 \olref[sth][choice][countablechoice]{sec}。} 如下定义单射 $g \colon
\bigcup_{\beta \in \cardfont{a}} X_\beta \to \cardfont{a} \times
\cardfont{a}$：$g(v) = \tuple{\beta, f_\beta(v)}$，其中 $v \in
X_\beta$ 且对任意 $\gamma \in \beta$ 有 $v \notin X_\gamma$。现在
$\bigcup_{\beta \in \cardfont{a}} X_\beta \preceq \cardfont{a} \times
\cardfont{a} \approx \cardfont{a}$ 由 \olref{cardplustimesmax} 即得。
\end{proof}

\end{document}